\section{Dense models and prediction}\label{sec:prediction}

We prove Principle~\ref{pr:prediction}.  The correlation estimate of
Section~\ref{sec:correlation} first allows the unbounded color weights to be
replaced by bounded functions.  We then choose nilsequence approximations at
coarse scales.  A Ramsey argument compares their projection energies with
those at the finer scales required by the correlation estimate.  All analytic
moduli in this last comparison will be independent of the number of master
indices.

Fix $m$.  Choose an integer $s\geq3$ such that
\begin{equation}\label{eq:prediction-step-choice}
 s+1\geq 2(2^d)-1
\end{equation}
for every cube order $d$ occurring in Proposition~\ref{prop:correlation-test}.
The order bound in that Proposition makes this a choice depending only on
$m$.  The later Cauchy--Schwarz argument will bound a $d$-cube test
by a subgroup norm of order $k=2^d$; concatenation then uses order
$2k-1$, whose inverse theorem gives
nilsequences of step $2k-2$.
Fix also the data $n,r,\chi,\mathcal B,\mathcal A$ in
Principle~\ref{pr:prediction}.  For the moment fix a master index count
$N$ and the parameters of Lemma~\ref{lem:master-scales}; only blocks
with at most $n$ members are used.
For $B=T\cup\{i\}$, $a\in\mathcal A$, and $c\in[r]$, write
\[
 \nu=\nu_B,\qquad f(y)=\1_{\chi(h_Bay)=c},\qquad \rho(y)=\nu(y)f(y).
\]
Here $f$ is zero on nonpositive integers.  All sufficiently large $w$ make
its positive color arguments integral.  We make two replacements,
$\rho\to F\to S$.  The bounded function
$F$ must approximate $\rho$ against the cube tests from
Section~\ref{sec:correlation};
the piecewise nilsequence $S$ must additionally admit alignment at the
coarse scale.  The gap-energy comparison will make these two demands
compatible.

\subsection{An algebra of dual tests}

A cube type at $B$ means one of the types supplied by
Proposition~\ref{prop:correlation-test}, at a gap valid for its chain.  We
also include, at every gap $\max T<l<i$, the type of dimension $s+1$ with
modulus $M$ and no prime variables.  Fix an integer $J_0\geq1$.  For any
one of these types put $M_{\boldsymbol p}=M(\boldsymbol p)$ and
\[
 L_{\boldsymbol p}=\floor{R_l/(J_0M_{\boldsymbol p})}.
\]
Primes have the good-tuple law of
Proposition~\ref{prop:correlation-test}; the variables $u_j^0,u_j^1$ are
independent and uniform in $[0,L_{\boldsymbol p})\cap\Z$.  Its dual tests
are the real functions
\begin{equation}\label{eq:prediction-dual-test}
 \mathcal D(y)=\E_{\boldsymbol p,u}e(\boldsymbol p)
 \prod_{\varnothing\ne\omega\subset[d]}
 g_{\omega,\boldsymbol p}
 \left(y+M_{\boldsymbol p}\sum_{j\in\omega}(u_j^1-u_j^0)\right),
 \qquad |e|\leq1,\quad |g_{\omega,\boldsymbol p}|\leq1+\nu.
\end{equation}
The functions, as well as the bounded prime factor, may vary arbitrarily
with $w$.  The finitely many cube templates and any fixed finite set of
$J_0$ values are included in each uniform assertion below.

To construct $F$, we need more than a mean-one estimate for $\nu$.
The next Lemma shows that $\nu-1$ is asymptotically orthogonal to the
algebra generated by the dual tests, and controls the tests sufficiently
to replace them by bounded ones.

\begin{lemma}[Dual-test pseudorandomness]\label{lem:dual-pseudorandomness}
For every fixed $b\geq0$ and dual tests
$\mathcal D_1,\ldots,\mathcal D_b$ at the same block,
\begin{equation}\label{eq:prediction-dual-products}
 \E_{\mu_i}(\nu-1)\prod_{q=1}^b\mathcal D_q=o(1)
\end{equation}
uniformly over their inputs.  The tests may use different permissible
gaps and different cube types.  If $d_*$ bounds the dimensions of these
types and $A_*=2^{2^{d_*}-1}$, then for every fixed positive integer $b$,
\begin{equation}\label{eq:prediction-dual-moments}
 \E_{\mu_i}(1+\nu)|\mathcal D|^b\leq2A_*^b+o(1).
\end{equation}
For any fixed $K>A_*$, replacing every test by its clipping to $[-K,K]$
changes it by $o(1)$ in every fixed $L^p((1+\nu)\mu_i)$ norm, and
\eqref{eq:prediction-dual-products} remains true for the clipped tests.
\end{lemma}

\begin{proof}
Expand the product using independent prime and shift replicas.  In replica
$q$, of dimension $d_q$ and modulus $M_q$, the nonroot rows are
\[
 Y_{q,\omega}=y+M_q\sum_{j\in\omega}(u_{q,j}^1-u_{q,j}^0),
 \qquad\varnothing\ne\omega\subset[d_q].
\]
The root row is $y$.  For a fixed nonempty $\omega$, choose integers
$a_0,a_1,\ldots,a_{d_q}$ such that
\[
 a_0\ne0,\qquad a_0+\sum_{j\in\omega}a_j=0,\qquad
 a_0+\sum_{j\in\omega'}a_j\ne0\quad(\omega'\ne\omega).
\]
Indeed, the hyperplane defined by the equality is not contained in any
of the finitely many hyperplanes excluded by the inequalities: their
normal vectors $(1,\1_{\omega'})$ are pairwise nonparallel.  A rational
point avoiding their intersections exists, and clearing denominators
gives the indicated integers.  Insert a translation parameter for this
row, acting by
\[
 y\longmapsto y+a_0M_qv,\qquad
 u_{q,j}^1\longmapsto u_{q,j}^1+a_jv,
\]
and leaving all other shift variables unchanged.  It fixes the chosen
row.  Every other row in its replica has a nonzero response, the root
has response $a_0M_q$, and a row in another replica has the same nonzero
response $a_0M_q$.

These directions are compatible even when the gaps differ.  Put
$V_B=2+M+\prod_{b\in T}X_b^2$.  The shortest shift interval in each
replica dominates every fixed power of $V_B$.  The translation parameters
for replica $q$ may, for example, have lengths comparable to the square
root of its shortest shift length.  Their lengths still dominate powers
of $V_B$, whereas the resulting displacements of its shifts are
negligible relative to their intervals, even after multiplication by any
fixed power of $V_B$.  The sum of the root displacements is likewise
negligible at $X_i$.  The parameters depend on the primes only through
the lengths, which is permitted by the uniform box version of
Proposition~\ref{prop:linear-forms}.  Thus inserting all these averaged
translations changes the original expression by $o(1)$: use the crude
bound $1+\nu\leq1+\prod_{b\in T}X_b^2$ and the translation estimates of
Lemma~\ref{lem:sampling}.  Fixed direction constants do not affect these
separations.

Eliminate the nonroot rows one at a time by the weighted
Cauchy--Schwarz procedure of Lemma~\ref{lem:additive-elimination}.
At the step assigned to a row, all its current copies are independent
of that row's translation parameter.  Bound those copies by their
$1+\nu$ weights and use their product as the outside density. Squaring
duplicates that row's assigned translation parameter; all other
variables remain in the outside expectation, as in that procedure.
Prime-only bounded factors can be discarded outside the
inner average.  Every prefactor is bounded, and one weight for each
eliminated copy is retained in the final expression.  The root, which
is never eliminated, contributes copies of $(\nu-1)$ indexed by all
choices of the duplicated translation parameters.

We verify the linear-forms hypotheses at every such application.
Every form has coefficient $1$ at $y$.  Different base rows have distinct
coefficients on the original shifts.  Copies of one base row are
distinguished by a duplicated direction with nonzero response; for an
eliminated row these are precisely directions other than its own.
Root copies are distinguished because every inserted direction has
nonzero root response.  A distinguishing minor is a fixed nonzero
integer times a single replica modulus $M_q$; no difference of two
moduli is required.  At a rough prime, its possible vanishing
is therefore included among the polynomial exceptional tests defining
that modulus.  Small-prime factors and fixed nonzero integer factors
are covered by the smooth modulus and, eventually, by $W$.
All replicas use gaps after the common tail $T$, so the multi-gap form
of Proposition~\ref{prop:linear-forms} applies with $V_B$.
The arbitrary prime-only restrictions in that Proposition include the
separate good-tuple restrictions in each replica.

Expand the retained $1+\nu$ weights and every root difference.  Each
resulting product has the main term obtained by replacing each $\nu$
by $1$, with a uniform $o(1)$ error.  In that substitution at least one
root difference becomes zero.  Thus the final Cauchy--Schwarz expression
is $o(1)$, proving \eqref{eq:prediction-dual-products}.  When $b=0$ this
is simply $\E\nu=1+o(1)$.

For the moment bound, Jensen's inequality and independent replicas give
\[
 |\mathcal D(y)|^b\leq
 \E\prod_{q=1}^b\prod_{\varnothing\ne\omega\subset[d]}
 (1+\nu)(Y_{q,\omega}).
\]
Include the root weight $1+\nu(y)$ and expand.  These rows are already
distinct without the additional translations.  Proposition~\ref{prop:linear-forms}
gives $2^{1+b(2^d-1)}+o(1)$, proving
\eqref{eq:prediction-dual-moments}.

Let $\mathcal D^{[K]}=\max(-K,\min(K,\mathcal D))$.  For integers $b>p$,
\[
 \limsup_{w\to\infty}
 \E(1+\nu)|\mathcal D-\mathcal D^{[K]}|^p
 \leq 2K^{p-b}A_*^b.
\]
This estimate holds for every fixed $b$; sending $b$ to infinity after
the limit proves that the left side is zero.  Telescoping a fixed
product and applying H\"older's inequality with respect to
$(1+\nu)\mu_i$ now transfers
\eqref{eq:prediction-dual-products} to clipped tests, since
$|\nu-1|\leq1+\nu$ and all required fixed moments are bounded.  No
moment order growing with $w$ has been used.
\end{proof}

We use the separation and polynomial-approximation proof of the dense
model theorem developed independently by Gowers
\cite[Section~4]{GowersDense} and Reingold, Trevisan, Tulsiani and
Vadhan \cite[Section~2]{RTTVDense}.
The preceding Lemma supplies the pseudorandomness needed for our
particular weighted test class.

\begin{proposition}[Bounded dense models]\label{prop:dense-model}
There are functions $F_{B,a,c}:\Z\to[0,1]$ such that, for every fixed
$J_0$ and every dual test at $B$,
\begin{equation}\label{eq:prediction-dense-approximation}
 \E_{\mu_i}(\rho-F_{B,a,c})\mathcal D=o(1)
\end{equation}
uniformly over its inputs and over the finitely many master blocks,
scale labels, and colors.
\end{proposition}

\begin{proof}
First fix a finite collection of $J_0$ values, a positive desired error,
and a common clipping bound $K$ from
Lemma~\ref{lem:dual-pseudorandomness}.  The support of $\mu_i$ is finite.
Let $\mathcal C$ be the closed convex hull of the signed clipped tests,
viewed as vectors on that support.  It is compact and lies in $[-K,K]$
coordinatewise.  The candidate models form the compact convex cube
$\mathcal F=[0,1]^{\supp\mu_i}$.  Finite-dimensional minimax for the
bilinear pairing gives
\[
 \inf_{F\in\mathcal F}\sup_{G\in\mathcal C}\langle\rho-F,G\rangle
 =\sup_{G\in\mathcal C}
 \bigl(\langle\rho,G\rangle-\langle1,G_+\rangle\bigr),
 \qquad G_+=\max(G,0),
\]
because $\sup_{F\in\mathcal F}\langle F,G\rangle=\langle1,G_+\rangle$.
Since $0\leq\rho\leq\nu$, the right-hand side is at most
$\sup_G\langle\nu-1,G_+\rangle$.

To bound it uniformly, approximate $t\mapsto\max(t,0)$ on $[-K,K]$
within $\delta$ by a polynomial $P(t)$.  Every fixed power of a convex
combination of signed tests is a convex combination of their signed
products.  Lemma~\ref{lem:dual-pseudorandomness}, followed by continuity
for the closed convex hull, therefore gives
$\langle\nu-1,P(G)\rangle=o(1)$ uniformly in $G\in\mathcal C$.  The
polynomial approximation costs at most
$\delta\E(1+\nu)=2\delta+o(1)$.  First take $w\to\infty$ and then
$\delta\to0$.  Minimax produces models with the desired error against
the clipped tests; the clipping estimate gives the same conclusion for
the original tests, since $|\rho-F|\leq1+\nu$.

Apply this argument successively to errors $1/q$ and $J_0\leq q$.
For each fixed $q$ it works for all sufficiently large $w$, uniformly
over the finite master data.  Choosing the stage $q=q(w)$ to increase
sufficiently slowly gives \eqref{eq:prediction-dense-approximation}
for every fixed $J_0$.  This diagonal choice is made together with the
fixed-template diagonal in Lemma~\ref{lem:master-scales}; it asserts no
estimate for an arbitrary growing moment order.  Extend the models
outside the support of $\mu_i$ with values in $[0,1]$.
\end{proof}

Here is the first consequence for a chain with a valid gap.  In its
weighted count, replace one factor
$\rho_{B_d,a_d,c}(L_J)$, $|J|\geq2$, by $F_{B_d,a_d,c}(L_J)$.
The difference has target $h=\rho-F$, with $|h|\leq1+\nu_B$;
all other linear factors still satisfy the bounds required by
Proposition~\ref{prop:correlation-test}.  The product masks and the
center weights are retained.  In the resulting cube, its nonroot
copies of $h$ are admissible inputs in
\eqref{eq:prediction-dual-test}; its root pairing is therefore $o(1)$
by \eqref{eq:prediction-dense-approximation}.  Equation~\ref{eq:correlation-test}
and a finite telescoping sum show that all nonsingleton color weights
can be replaced by their bounded models with total error $o(1)$.

\subsection{Nilsequence tests and nested projections}

The bounded functions $F$ now replace the weighted color factors in the
counts.  Calibration requires another comparison: $\rho$ and $F$ must
agree against bounded Lipschitz functions of the eventual prediction
$S$, so that small prediction values cannot carry much actual color
mass.  We first establish this nilsequence testing property, using
missing-corner reconstruction to express the tests through the dual
tests already controlled.  We then choose coarse models whose
differences from $F$ have small projections at the finer chain scales.
The next subsection turns those projection bounds into counting estimates.

\begin{lemma}[Testing piecewise nilsequences]\label{lem:nilsequence-testing}
Fix $B=T\cup\{i\}$ and a gap $\max T<l<i$.  Let $S'$ be any family
of bounded real piecewise nilsequences on intervals of length $R_l$
and residues modulo $M$, of step at most $s$ and bounded complexity
in the sense of Definition~\ref{def:piecewise-model}.  Then
\begin{equation}\label{eq:prediction-nilsequence-testing}
 \E_{\mu_i}(\rho-F_{B,a,c})S'=o(1).
\end{equation}
The assertion is uniform for each fixed finite list of nilmanifolds
and fixed observable bounds.
\end{lemma}

\begin{proof}
Lemma~\ref{lem:cube-corner}, proved independently in
Section~\ref{sec:cube-limits}, supplies a compact set of nilmanifold
cubes containing every linear orbit cube
\[
 (x_\omega)_{\omega\subset[s+1]}
   =\left(g^{k+\sum_{j\in\omega}v_j}x\right)_{\omega\subset[s+1]},
 \qquad k,v_1,\ldots,v_{s+1}\in\Z.
\]
On this set the nonroot vertices determine $x_\varnothing$
continuously.  Its Stone--Weierstrass consequence gives the following
concrete input.  For any accuracy $\delta>0$, the observable at the
root is uniformly approximated by a finite sum
\[
 \sum_{t=1}^{q}\lambda_t
 \prod_{\varnothing\ne\omega\subset[s+1]}\phi_{t,\omega}(x_\omega)
\]
with $|\phi_{t,\omega}|\leq1$.  The same Lemma gives finitely many
such recipes, with fixed total coefficient bounds, for our finite
nilmanifold list and uniformly bounded Lipschitz observables.
The recipes are valid for every $g$ and $x$, so they can be used on
all pieces even though the translating elements and base points are
unrestricted.  It remains to realize these recipes as our dual tests
and control the cubes that cross piece boundaries.

Use the extra dual test of dimension $s+1$, modulus $M$, and a fixed
large $J_0$.  If $y$ is farther than $(s+1)R_l/J_0$ from the endpoints
of its $R_l$ interval, all its cube vertices lie in that interval and
have the same residue modulo $M$.  The recipe for that piece then
reconstructs $S'(y)$.  Recipes can be selected through the input
functions of the dual test: for a term belonging to recipe $r$, put
in every nonroot input the indicator that its piece selected $r$,
multiplied by its single-vertex factor evaluated on that piece's
orbit.  Sum over the finite recipes and terms.  On cubes staying in
one piece this selects exactly its recipe.  Each resulting input is
bounded by $1$, and hence is allowed in
\eqref{eq:prediction-dual-test}.

We record the weighted boundary estimate needed for the discarded
roots.  If $E_l$ is the union of these boundary strips, then
\begin{equation}\label{eq:prediction-weighted-boundary}
 \E_{\mu_i}(1+\nu_B)\1_{E_l}
 \leq O_s(J_0^{-1})+o(1).
\end{equation}
For a fixed divisor draw $\sigma=t_T$, divisibility by $\sigma$
and coprimality to $W$ are periodic with period $\sigma W$.
This period is negligible compared with $R_l$, uniformly in
$\sigma\leq\prod_{b\in T}X_b^2$.  Counting in each full $R_l$ interval
shows that a union of end strips of relative length $O_s(J_0^{-1})$
has that proportion of its weighted mass, with relative error
$O(\sigma W/R_l)$.  The harmonic factor varies by $o(1)$ within a
cell, since $R_l/X_i=o(1)$; the two partial cutoff cells have $o(1)$
mass by the same counting estimate.  This proves the assertion for
$\sigma\1_{\sigma\mid y}\mu_i$, uniformly in $\sigma$.
Average over $\sigma$ and also take $\sigma=1$ to obtain
\eqref{eq:prediction-weighted-boundary}.

The recipe approximation costs $O(\delta)$ against
$|\rho-F|\mu_i\leq(1+\nu_B)\mu_i$.  For fixed recipes and $J_0$,
Proposition~\ref{prop:dense-model} makes every dual pairing $o(1)$.
The terms outside the good roots are bounded by a fixed constant
(depending on the recipes) times the quantity in
\eqref{eq:prediction-weighted-boundary}.  First fix sufficiently
accurate recipes, then take $J_0$ sufficiently large for those
recipes, and finally take $w\to\infty$.  Since $\delta$ was arbitrary,
this proves \eqref{eq:prediction-nilsequence-testing}.
\end{proof}

Use the fixed nonprincipal ultrafilter $\mathcal U$ on the positive
integers $w$. For each pivot $i$, form the real Hilbert space of families of
vectors in $L^2(\mu_i)$ with uniformly bounded norms, with pairing
$\lim_{\mathcal U}\langle\cdot,\cdot\rangle$, after quotienting out
zero-norm families and completing.  For $l<i$, let $\mathcal N_{i,l}$
be the closed span in this space of bounded-complexity piecewise
step-at-most-$s$ nilsequence families on the $R_l$ intervals and
residues modulo $M$.  Complexity is allowed to vary from one family
to another but is bounded within a family.  Denote the orthogonal
projection by $P_{i,l}$.

When $l<l'<i$, the divisibility $R_l\mid R_{l'}$ aligns the smaller
intervals with the larger ones.  Restricting an orbit to a smaller
interval only changes its base point.  Consequently
\begin{equation}\label{eq:prediction-nested-spaces}
 \mathcal N_{i,l'}\subseteq\mathcal N_{i,l},\qquad
 \|P_{i,l}v-P_{i,l'}v\|_2^2
 =\|P_{i,l}v\|_2^2-\|P_{i,l'}v\|_2^2.
\end{equation}
Products and Lipschitz real combinations of finitely many representing
families remain representing families, using product nilmanifolds.
This includes clipping to $[0,1]$.

\begin{lemma}[Ramsey selection of gap energies]\label{lem:energy-selection}
For every $\eps>0$, a master count $N$ depending only on
$n,r,|\mathcal A|,\eps$ can be chosen so that there are indices
\[
 k_1<j_1<k_2<j_2<\cdots<k_n<j_n
\]
and $[0,1]$-valued piecewise models $S_{B,a,c}$ at pivot $i=j_u$
on intervals of length $R_{k_u}$ such that
\begin{equation}\label{eq:prediction-coarse-approximation}
 \|S_{B,a,c}-P_{j_u,k_u}F_{B,a,c}\|_2\leq\eps.
\end{equation}
For a chain on the principal indices $j_1,\ldots,j_n$, let $l$ be the
padding index immediately before its first pivot.  At every block
$B=T\cup\{j_u\}$ of the chain,
\begin{equation}\label{eq:prediction-fine-projection}
 \|P_{j_u,l}(F_{B,a,c}-S_{B,a,c})\|_2\leq2\eps.
\end{equation}
The families $S$ have bounded complexity after $N$ and $\eps$ are fixed.
\end{lemma}

\begin{proof}
The energies $\|P_{i,l}F_{T\cup\{i\},a,c}\|_2^2$ lie in $[0,1]$.
Color each ordered tuple consisting of the increasing members of
$T$, followed by $l$ and $i$, by their bins of width $\eps^2$,
jointly for all $a,c$.  There are finitely many colors, with a bound
independent of $N$.  Apply finite Ramsey successively for all tuple
lengths $3,\ldots,n+1$ to obtain a homogeneous set of size $2n$.
Label it alternately by padding and principal indices as displayed.

At each selected block approximate its coarse projection by a finite
linear combination $T'$ of representing families, to accuracy $\eps$.
Its clipping $S$ to $[0,1]$ is still a representing family.  Pointwise
clipping decreases the distance to $F\in[0,1]$.  Since both $T'$ and
$S$ lie in the coarse subspace, orthogonality gives
\[
 \|F-S\|_2^2=\|F-PF\|_2^2+\|PF-S\|_2^2,
\]
and the analogous identity for $T'$.  Thus clipping also decreases
the distance to $PF$, proving
\eqref{eq:prediction-coarse-approximation}.  There are only finitely
many selected blocks and labels, so their representing families
share a finite nilmanifold list and uniform observable bounds.

All tails of the chain precede its first pivot and hence precede
$l$; also $l\leq k_u<j_u$.  If $l<k_u$, the tuples $(T,l,j_u)$ and
$(T,k_u,j_u)$ lie in the homogeneous set and have the same color.
Their energies differ by at most $\eps^2$.
Equation~\eqref{eq:prediction-nested-spaces} bounds the distance
between the two projections by $\eps$.  As $S$ belongs to the coarse
and hence the fine space, \eqref{eq:prediction-fine-projection}
follows by the triangle inequality.  The case $l=k_u$ follows
directly from \eqref{eq:prediction-coarse-approximation}.
\end{proof}

\subsection{Subgroup cubes and the inverse theorem}

The energy selection makes $F-S$ have small projection onto the fine
nilsequence space.  We must turn that conclusion into a small cube test
in Proposition~\ref{prop:correlation-test}, with a threshold independent
of the master count $N$.  Finite subgroup norms supply this link.

For a finite subgroup $Q$ acting by measure-preserving transformations
$T^q$ on a probability space $\mathcal X$, and a real bounded function
$h$, use the convention
\[
 \|h\|_{U^t_Q(\mathcal X)}^{2^t}
 =\E_{x\in\mathcal X}\E_{v_1,\ldots,v_t\in Q}
 \prod_{\omega\subset[t]}h\left(T^{\sum_{j\in\omega}v_j}x\right).
\]
This is the usual subgroup Gowers seminorm; for complex functions the
usual alternating conjugations are inserted.

We use exactly two external inverse statements.  First, the subgroup
case of the finitary qualitative Bessel inequality
\cite[Theorem~1.23]{TaoZiegler} says that for each $k$ there is a
function $b_k(\delta)\to0$ as $\delta\downarrow0$ such that
\begin{equation}\label{eq:prediction-bessel}
 \E_{\alpha,\alpha'}
 \|h\|_{U^{2k-1}_{Q_\alpha+Q_{\alpha'}}(\mathcal X)}\leq\delta
 \quad\Longrightarrow\quad
 \E_\alpha\|h\|_{U^k_{Q_\alpha}(\mathcal X)}\leq b_k(\delta),
 \qquad |h|\leq1.
\end{equation}
The bound is independent of the finite family, the group, and the
probability system.  Finite subgroups are rank-zero coset
progressions in that Theorem.  Its dilations leave them unchanged,
and its multiset sum of two subgroups induces uniform measure on
their subgroup sum.  The statement consequently has exactly the
form \eqref{eq:prediction-bessel}.  Rational probability weights on
indices are implemented by repetition of indices; arbitrary finite
weights follow by approximation.  We choose $b_k$ nondecreasing and
enlarge its argument by a factor of two to absorb the approximation
slack; the resulting modulus, still denoted $b_k$, tends to zero.

Second, the Green--Tao--Ziegler inverse theorem, in its finite-list
linear-nilsequence formulation
\cite[Conjecture~1.2 and Theorem~1.3]{GTZ}, with the correction
\cite{GTZErratum}, gives the following statement.  For each integer
$t\geq2$ and $\delta>0$, a $1$-bounded function on an integer interval
with $U^t$ norm at least $\delta$ correlates by at least
$c_t(\delta)>0$ with a sequence $\Phi(g^kx)$ on a nilmanifold of
step at most $t-1$.  There is a finite list of such nilmanifolds,
with connected simply connected groups, from which the manifold is
chosen; $\|\Phi\|_\infty\leq1$ and its Lipschitz bound depend only on
$t,\delta$.  The cases $t=2,3$ are the established lower-order
cases recalled in the Introduction of \cite{GTZ}; Theorem~1.3 supplies
the higher orders.  No bound on $g$ is asserted or needed.
We use only this inverse conclusion, not the original auxiliary
Proposition~8.3 addressed by the erratum.  The interval norm here is
the normalized cube mean with all vertices in the interval, equivalently
the zero-extension norm divided by the norm of the interval indicator
in a sufficiently large auxiliary cyclic group.

\begin{lemma}[From subgroup cubes to a fine nilsequence projection]
\label{lem:subgroup-inverse}
Fix a cube type of dimension $d$, a gap $l<i$, and $J_0\geq1$.
Assume $2(2^d)-2\leq s$.  For every $\gamma>0$ there is
$\kappa=\kappa(d,J_0,\gamma)>0$, independent of the master count,
the gap, the cutoffs, and the cell probabilities, with the following
property.  If $h:\Z\to[-1,1]$ is a family with
$\|P_{i,l}h\|_2<\kappa$, then
\begin{equation}\label{eq:prediction-subgroup-conclusion}
 \lim_{\mathcal U}\left|
 \E_{\boldsymbol p,y,u}
 \prod_{\omega\subset[d]}
 h\left(y+M_{\boldsymbol p}
       \sum_{j\in\omega}(u_j^1-u_j^0)\right)\right|
 \leq\gamma+O_d(J_0^{-1}).
\end{equation}
Here $y\sim\mu_i$, and primes and shifts have the laws in
\eqref{eq:prediction-dual-test}.  The constant in the last term
depends only on $d$.
\end{lemma}

\begin{proof}
We pass from the short increments to subgroup norms, combine the
subgroups using rough coprimality, and apply the inverse theorem in
each cyclic cell.  We keep the cell probabilities throughout, so every
threshold is independent of the number of cells and the master count.

\paragraph{From short increments to subgroup norms.}
Discard the two partial $R_l$ intervals at the ends of
$[X_i,X_i^2)$; their $\mu_i$ mass is $o(1)$.  Within every remaining
interval and every residue modulo $M$, harmonic measure differs
from conditional uniform measure by relative $o(1)$, uniformly in
the cell, because $R_l/X_i=o(1)$.  Only residues coprime to $W$
occur, and translation by $M$ preserves them.

Write $q_l=R_l/M$.  Each cell has exactly $q_l$ points.  Give it its
original probability, normalized after discarding the partial
cells, and identify its progression indices with
$G_l=\Z/q_l\Z$.  Their disjoint union, with these probabilities and
conditional Haar measure in each cell, is a finite probability
system $\mathcal X_l$.  Translation by $1\in G_l$ means translation
by $M$ with wrapping inside that cell.  Replacing the original cube
by this periodized cube costs $O_d(J_0^{-1})+o(1)$: every vertex
moves by at most $dR_l/J_0$ from its root, so only roots within that
distance of an interval endpoint can wrap.  All functions in this
argument are bounded by $1$.

For a prime tuple put
\[
 q_{\boldsymbol p}=M_{\boldsymbol p}/M,
 \qquad Q_{\boldsymbol p}=q_{\boldsymbol p}G_l.
\]
The divisibility provisions in Lemma~\ref{lem:master-scales} give
$q_{\boldsymbol p}\mid q_l$.  Its subgroup has order
$q_l/q_{\boldsymbol p}$, while
$L_{\boldsymbol p}=\floor{q_l/(J_0q_{\boldsymbol p})}$ tends to
infinity uniformly over the pool.  The pushforward of a uniform
$u\in[0,L_{\boldsymbol p})$ to $q_{\boldsymbol p}u\in Q_{\boldsymbol p}$
has density at most $2J_0$ relative to subgroup Haar measure.
The difference of two independent such variables has no larger
density.  Thus the $d$ independent increments of the periodized
cube have joint density at most $(2J_0)^d$ on
$Q_{\boldsymbol p}^d$.

For completeness, this density can be removed with a bounded number
of Cauchy--Schwarz steps.  Regard it as one factor, independent of
the root, on Haar variables $(x,v_1,\ldots,v_d)$, where
$x\in\mathcal X_l$ and $v_j\in Q_{\boldsymbol p}$.
Insert one uniform subgroup direction translating the root alone;
the density is invariant in that direction.  For each nonempty
$\omega\subset[d]$, choose $j\in\omega$ and insert a direction
\[
 x\longmapsto T^a x,\qquad v_j\longmapsto v_j-a,
 \qquad a\in Q_{\boldsymbol p},
\]
which leaves that vertex fixed.  These invariant changes of Haar
variables commute.  Successive Cauchy--Schwarz inequalities outside
the assigned direction discard first the density and then each
nonroot vertex function, including its copies from earlier steps.
The root has response $a$ in every direction.  After
$k=1+(2^d-1)=2^d$ steps its retained copies are a Haar subgroup
$k$-cube.  Taking the $2^k$-th root proves
\begin{equation}\label{eq:prediction-cube-subgroup-bound}
 |\text{periodized cube at }\boldsymbol p|
 \leq C_{d,J_0}\|h\|_{U^k_{Q_{\boldsymbol p}}(\mathcal X_l)}.
\end{equation}
The argument acts inside each cell and integrates its probability
throughout; its constant is independent of the number of cells.

\paragraph{Combining the subgroups.}
For independent tuples, Lemma~\ref{lem:rough-coprimality}, also after
the good-tuple restrictions, gives
$\gcd(q_{\boldsymbol p},q_{\boldsymbol p'})=1$ with probability
$1-o(1)$.  On this event B\'ezout's identity implies
$Q_{\boldsymbol p}+Q_{\boldsymbol p'}=G_l$.  Consequently
\[
 \E_{\boldsymbol p,\boldsymbol p'}
 \|h\|_{U^{2k-1}_{Q_{\boldsymbol p}+Q_{\boldsymbol p'}}(\mathcal X_l)}
 \leq \|h\|_{U^{2k-1}_{G_l}(\mathcal X_l)}+o(1).
\]
Together with \eqref{eq:prediction-bessel} and
\eqref{eq:prediction-cube-subgroup-bound}, this says: for every
$\gamma>0$ a bound on the global $U^t$ norm, $t=2k-1$, sufficiently
small in terms of $d,J_0,\gamma$ makes the periodized cube mean at
most $\gamma$.  This assertion is uniform in all finite system data.

\paragraph{From the global norm to a nilsequence projection.}
We next show that a fixed positive global norm forces a fixed
positive fine projection.  We first justify use of the interval
inverse theorem on one cyclic cell, without requiring the resulting
nilsequence to be periodic.  Let $v$ be $1$-bounded on $\Z/q\Z$ and
extend it periodically to $[0,Kq)$.  For $L=Kq$ define the integer
cube domain
\[
 A_L=\{(x,a_1,\ldots,a_t)\in\Z^{t+1}:
               0\leq x+\textstyle\sum_{j\in\omega}a_j<L
               \text{ for every }\omega\subset[t]\}.
\]
The $2^t$-th power of the interval norm is the mean of the periodic
cube product over $A_L$.  Averaging a parameter translate
$r\in\{0,\ldots,q-1\}^{t+1}$ gives exactly the cyclic cube mean
at every parameter point.  The discrepancy is therefore bounded
by $\E_r|(A_L+r)\mathbin\triangle A_L|/|A_L|$.
For a vertex $\omega$, its displacement is
$r_0+\sum_{j\in\omega}r_j\leq(t+1)(q-1)$.
The boundary discrepancy at that vertex has size at most this
quantity times $2L^t$: once that vertex and its $t$ neighboring
vertices are specified the integer cube is determined, and each
neighbor has at most $L$ choices on the relevant side of the
symmetric difference.  Summing over vertices bounds the numerator
by $O_t(qL^t)$.  Nonnegative increments and root with total less
than $L$ already give
$|A_L|\geq\binom{L+t}{t+1}\geq L^{t+1}/(t+1)!$.
Thus, uniformly in $q,K,v$,
\begin{equation}\label{eq:prediction-cyclic-interval}
 \|v^{\mathrm{per}}\|_{U^t[Kq]}^{2^t}
 =\|v\|_{U^t(\Z/q\Z)}^{2^t}+O_t(K^{-1}).
\end{equation}
This is an estimate for the powers of the norms.

If the cyclic norm is at least $\delta$, choose a sufficiently
large fixed $K=K(t,\delta)$ in
\eqref{eq:prediction-cyclic-interval}.  The interval norm is then
at least $\delta/2$.  The inverse theorem gives a uniformly bounded
complexity correlator on those $K$ periods.  Splitting the
correlation into its $K$ period averages shows that one period has
correlation of at least the same modulus.  Translating that period
back to $[0,q)$ only replaces the nilsequence base point by a
translate $g^{jq}x$.  It leaves its manifold, step and observable
bounds unchanged.  We have proved a cyclic-to-interval correlation
statement on fixed representatives of $\Z/q\Z$; the correlator
itself need not have period $q$.

Now write the cell weights of $\mathcal X_l$ as $\alpha_C$.
The exact identity
\[
 \|h\|_{U^t_{G_l}(\mathcal X_l)}^{2^t}
 =\sum_C\alpha_C\|h_C\|_{U^t(G_l)}^{2^t}
\]
shows that, if the left-hand norm is at least $\delta$, cells with
$\|h_C\|_{U^t(G_l)}\geq\delta/2$ have total probability at least
$\delta^{2^t}/2$.  On each such cell choose its inverse-theorem
correlator, rotate its complex phase and take the real part so
that its pairing with the real function $h_C$ is positive.
Put zero on the remaining cells.  Every chosen sequence comes
from the same finite nilmanifold list and has the same uniform
observable bounds.  They form one permitted real piecewise
nilsequence family $V$ at gap $l$, with $|V|\leq1$.  Translating
from conditional uniform measure back to $\mu_i$ costs $o(1)$.
Thus along any set of $w$ on which the global norm is at~least~$\delta$,
\[
 \langle h,V\rangle\geq
 \tfrac12\delta^{2^t}c_t(\delta/4)+o(1).
\]
The harmless smaller inverse threshold allows slack in all these
inequalities.  The step is at most $t-1=2(2^d)-2\leq s$, so
$V$ represents an element of $\mathcal N_{i,l}$.  If the stated
norm bound holds on a set belonging to $\mathcal U$, select the
correlators there and set the family to zero on its complement.
Orthogonality and $\|V\|_2\leq1$ give the positive lower bound
\[
 \|P_{i,l}h\|_2\geq
 \tfrac12\delta^{2^t}c_t(\delta/4).
\]

Choose $\delta$ small enough that Bessel and
\eqref{eq:prediction-cube-subgroup-bound} give the prescribed
$\gamma$, and choose $\kappa$ smaller than the displayed positive
projection bound.  If the projection is less than $\kappa$, the
global norm cannot exceed that threshold along $\mathcal U$.
Restoring the periodization error proves
\eqref{eq:prediction-subgroup-conclusion}.
\end{proof}

\subsection{Completion of the Prediction Principle}

\begin{proof}[Proof of Principle~\ref{pr:prediction}]
Fix the data $n,r,\chi,\mathcal B,\mathcal A,\tau,\eta$ of the
Principle, with $s=s(m)$ as in
\eqref{eq:prediction-step-choice}.  There are
$q_m=2^m-m-1$ nonsingleton factors in each count.  We spell out the
order of choices before selecting master indices.

Choose a positive target cube error $\zeta$ so small that
\[
 C_m(2\zeta)^\theta<\eta/(2q_m)
\]
for every exponent and constant in
Proposition~\ref{prop:correlation-test}.  There are only finitely
many orders and exponent choices, all controlled by $m$; increasing
the uniform constant if necessary makes this one finite set of
requirements.  Choose a fixed $J_0$ so large that every
periodization term $O_d(J_0^{-1})$ in
Lemma~\ref{lem:subgroup-inverse} is smaller than $\zeta$.
For the finitely many correlation-test types, apply that Lemma with
$\gamma=\zeta$, and choose
\[
 0<\eps<\eta,
 \qquad 2\eps<\min_d\kappa(d,J_0,\zeta).
\]
These choices depend on the requested tolerances and on $m$,
and are independent of $N$.  Now choose $N$ by
Lemma~\ref{lem:energy-selection}, construct its arithmetic scales
and dense models, and carry out that Lemma.  Restrict $h_i,X_i,t_i$
to the principal indices $j_1,\ldots,j_n$, relabel them by $[n]$,
and set $H_u=R_{k_u}$.  The intervening padding gaps and
Lemma~\ref{lem:master-scales} give all conditions in
Definition~\ref{def:admissible}; the models $S$ have exactly the
pieces required by Definition~\ref{def:piecewise-model}.

For calibration at a selected block choose a Lipschitz function
$\psi:[0,1]\to[0,1]$ equal to $1$ on $[0,2\tau]$ and to $0$ on
$[3\tau,1]$.  The bounded distribution replacement in
Corollary~\ref{cor:product-law} gives
\[
 \Pr\bigl(\chi(h_Bat_B)=c,\ S_{B,a,c}(t_B)\leq2\tau\bigr)
 \leq\E_{\mu_i}\rho\,\psi(S_{B,a,c})+o(1).
\]
The function $\psi(S)$ is a bounded-complexity coarse piecewise
nilsequence.  Lemma~\ref{lem:nilsequence-testing} replaces $\rho$
by $F$.  In the Hilbert limit it also lies in the coarse subspace,
so, writing $P=P_{i,k_u}$ and using
\eqref{eq:prediction-coarse-approximation},
\[
 \langle F,\psi(S)\rangle
 =\langle PF,\psi(S)\rangle
 \leq\langle S,\psi(S)\rangle+\eps
 \leq3\tau+\eps.
\]
This is the first assertion of the Principle, since $\eps<\eta$.

For the count comparison, fix a chain on the principal indices,
$b\in\mathcal B$, and a color.  Its block scale $b_{B_d}$ belongs
to $\mathcal A$, so all its models have already been constructed.
As shown after Proposition~\ref{prop:dense-model}, telescope the
nonsingleton color weights $\rho_d(L_J)$ to $F_d(L_J)$ at error
$o(1)$, retaining the center weights $\nu_d(z_d)$ and every product
mask, including the singleton masks.

Next telescope $F_d(L_J)$ to $S_d(L_J)$.  Choose as fine gap the
padding index immediately before the chain's first pivot.  It is
valid for every block of the chain.  Each individual difference
has bounded target $h=F_d-S_d\in[-1,1]$, and
\eqref{eq:prediction-fine-projection} gives
$\|P_{i_d,l}h\|_2\leq2\eps$.  All other factors, including the
unbounded center weights, still satisfy the bounds of
Proposition~\ref{prop:correlation-test}.  The parameter choices and
Lemma~\ref{lem:subgroup-inverse} bound its target cube mean by
$2\zeta$ in the ultrafilter limit.  Equation~\ref{eq:correlation-test}
therefore bounds each telescoping error by less than
$\eta/(2q_m)$.  Summing the finitely many differences gives a
limit error at most $\eta/2$ between the original count and
\[
 \E_{z\sim\bigotimes_d\mu_{i_d}}
 U(z)\prod_d\nu_d(z_d)
       \prod_{|J|\geq2}S_{d(J)}(L_J(z)).
\]
All factors except the center weights are now bounded by $1$.
The blocks of the chain are disjoint.  Corollary~\ref{cor:product-law}
thus replaces this last expectation, with error $o(1)$, by
\[
 \E_t U\bigl(z(t)\bigr)
       \prod_{|J|\geq2}S_{d(J)}\bigl(L_J(z(t))\bigr),
 \qquad z(t)_d=t_{B_d}.
\]
This proves the required count comparison with the stated tolerance
$\eta$.  Every limit in this construction may be taken along the
same fixed ultrafilter $\mathcal U$, as allowed in
Principle~\ref{pr:prediction}.  In particular the auxiliary
$\eps$ was chosen before the master Ramsey count $N$, whereas the
resulting model complexity was allowed to depend on all those
fixed choices.  The proof of the Principle is complete.
\end{proof}
